\section{Tensor products and stable finiteness}
\label{sec:consequences}

The consequences in Corollary~\ref{cor:stable-finiteness} use the
established factor form of the quasitrace problem and a trace-obstruction
theorem for tensor products. We record their scopes before applying them
to the different algebras in the corollary.

\paragraph{The factor input.}
The assertion that every normalized \(2\)-quasitrace on every unital
\(C^*\)-algebra is a tracial state is equivalent to the assertion that
every \(AW^*\)-factor of type \(\mathrm{II}_1\) is a \(W^*\)-algebra;
see \cite[Section~2]{MR} and \cite[Theorem~4.11]{Gow}.
Every such factor admits
a normalized \(2\)-quasitrace, and it admits a tracial state exactly
when it is a \(W^*\)-algebra \cite[Section~2]{MR}. A finite
\(AW^*\)-factor is simple as a \(C^*\)-algebra
\cite[Remark~2.7]{Gow}. For a simple unital \(C^*\)-algebra, absence of
stable proper infiniteness is equivalent to stable finiteness
\cite[Section~2]{MR}. These statements impose no separability condition
on the factor.

\paragraph{The tensor input and the passage to \(\mathbb F_2\).}
Haagerup's trace-obstruction theorem, in the properly infinite
formulation recorded by Milh{\o}j and R{\o}rdam
\cite[Theorem~3.1]{MR}, says that a unital \(C^*\)-algebra \(D\) has no
tracial state if and only if
\(D\otimes_{\min}C_r^*(\mathbb F_\infty)\) is properly infinite.
Haagerup notes that the subgroup inclusion
\(\mathbb F_\infty\hookrightarrow\mathbb F_2\) induces a unital
inclusion \(C_r^*(\mathbb F_\infty)\hookrightarrow C_r^*(\mathbb F_2)\)
\cite[Remark~2.5(a)]{Haagerup}. The minimal tensor product preserves
injective \(*\)-homomorphisms, so it gives a unital inclusion
\[
D\otimes_{\min}C_r^*(\mathbb F_\infty)
\hookrightarrow D\otimes_{\min}C_r^*(\mathbb F_2).
\]
The two isometries witnessing proper infiniteness retain their relations
under this inclusion. Consequently,
\begin{equation}
 D\text{ has no tracial state}
 \quad\Longrightarrow\quad
 D\otimes_{\min}C_r^*(\mathbb F_2)\text{ is properly infinite}.
 \label{eq:traceless-tensor}
\end{equation}

\begin{proof}[Proof of Corollary~\ref{cor:stable-finiteness}]
For (i), Theorem~\ref{thm:main} and the factor input give an
\(AW^*\)-factor \(M\) of type \(\mathrm{II}_1\) that is not a
\(W^*\)-algebra. Hence \(M\) has no tracial state. It is simple and
admits a normalized \(2\)-quasitrace. Theorem~\ref{thm:existence}
therefore excludes stable proper infiniteness, and the criterion for
simple algebras makes \(M\) stably finite.

Powers's theorem \cite{Powers} gives simplicity of
\(R=C_r^*(\mathbb F_2)\). Its canonical tracial state is faithful, as
are the matrix traces induced by it. A faithful trace rules out a
nonunitary isometry, so \(R\) is stably finite. It is also unital.
Applying \eqref{eq:traceless-tensor} to \(M\) proves (i):
\(M\otimes_{\min}R\) is properly infinite.

For (ii), take the explicit algebra \(B\) of
Theorem~\ref{thm:construction}. No matrix algebra over \(B\) is properly
infinite, so Theorem~\ref{thm:existence} supplies a normalized
\(2\)-quasitrace. The relations \eqref{eq:algebra-relations} exclude a
tracial state: applying one to their two sums would give \(1\le2/3\).
Equation~\eqref{eq:traceless-tensor} makes
\(B\otimes_{\min}R\) properly infinite, and
Theorem~\ref{thm:existence} then excludes a normalized \(2\)-quasitrace
on that tensor product.

For (iii), Haagerup's finite trace-obstruction criterion
\cite[Lemma~2.1]{Haagerup} supplies \(y_1,\ldots,y_n\in M\) such that
\[
 \sum_{i=1}^n y_i^*y_i=1,
 \qquad
 \left\|\sum_{i=1}^n y_iy_i^*\right\|<1.
\]
Let \(C=C^*(1,y_1,\ldots,y_n)\subset M\). This algebra is separable and
unital. It inherits stable finiteness from \(M\): an isometry in
\(M_k(C)\) is an isometry in \(M_k(M)\), hence is unitary there and
therefore already unitary in \(M_k(C)\). A tracial state \(\tau\) on
\(C\) would satisfy
\[
 1=\tau\left(\sum_{i=1}^n y_i^*y_i\right)
  =\tau\left(\sum_{i=1}^n y_iy_i^*\right)
  \le\left\|\sum_{i=1}^n y_iy_i^*\right\|<1,
\]
a contradiction. This proves (iii).
\end{proof}
